\section{Selection}
\label{sec:selection}

\subsection{Profile}

Table~\ref{tab:matrix} scores the profiles at $L/D = 4.5$ on the criteria of
Section~\ref{sec:options}. Heating is left out because, with a metallic tip, it does not depend
on the profile (Section~\ref{sec:heating}).

\begin{table}[H]
\centering
\caption{Trade-off matrix at $L/D = 4.5$ (1 = worst, 5 = best). Performance: apogee at equal length (Table~\ref{tab:trade}). Integration: 3U length inside the curved part. Manufacture: tangent shoulder, mould and tip. Heritage: flown or built at ERT. Stability: change in static margin.}
\label{tab:matrix}
\small
\begin{adjustbox}{max width=\textwidth}
\begin{tabular}{@{}lccccccc@{}}
\toprule
Criterion (weight) & Cone & Ogive & Parab. & Power $\tfrac12$ & Power $\tfrac34$ & LV-Haack & Von K\'arm\'an \\
\midrule
Performance (0.40) & 1.0 & 2.6 & 2.6 & 4.3 & 4.4 & 2.9 & 5.0 \\
Integration, 3U fit (0.20) & 1.0 & 4.8 & 4.8 & 3.0 & 1.7 & 5.0 & 3.7 \\
Manufacture (0.20) & 3.0 & 5.0 & 5.0 & 3.0 & 2.0 & 5.0 & 5.0 \\
Heritage (0.10) & 1.0 & 4.0 & 2.0 & 5.0 & 1.0 & 1.0 & 2.0 \\
Stability (0.10) & 5.0 & 1.4 & 1.5 & 2.2 & 3.9 & 1.0 & 2.1 \\
\midrule
\textbf{Weighted total} & 1.80 & 3.54 & 3.36 & 3.64 & 3.00 & 3.37 & \textbf{4.15} \\
\bottomrule
\end{tabular}
\end{adjustbox}
\end{table}


\textbf{Von K\'arm\'an is selected.} Three profiles form a leading group on flight performance:
Von K\'arm\'an, power $\tfrac12$ and power $\tfrac34$. At the design length, Von K\'arm\'an is
first. It has the lowest wave drag at every computed Mach number, flies \SI{18}{m} higher than
power $\tfrac34$ and \SI{22}{m} higher than power $\tfrac12$, and stays first in all seven robustness
scenarios (Figure~\ref{fig:robustness}). Ogive, parabolic and LV-Haack lose \SIrange{69}{79}{m},
and the cone \SI{132}{m}. Von K\'arm\'an also wins on the other criteria:
\begin{itemize}[leftmargin=*,itemsep=2pt]
\item it is tangent to the body at the shoulder, so there is no kink at the coupler bond. The
power-$\tfrac12$ nose has a $3.2^\circ$ kink and the power-$\tfrac34$ a $4.8^\circ$ kink at
$L/D=4.5$;
\item it holds more of the 3U unit (\FINTHREEU{} inside the curved part, against
\SI{217}{mm} for power $\tfrac12$ and \SI{152}{mm} for power $\tfrac34$);
\item its advantage holds at $M \geq 1.2$, so it does not hinge on the transonic model.
\end{itemize}
The Firehorn power-$\tfrac12$ nose remains an acceptable alternative (\SI{22}{m} lower) if
reusing its tooling outweighs the shoulder kink (R10, Section~\ref{sec:recommendations}).

\subsection{Length}

Every \SI{90}{mm} of nose ($\Delta L/D = 0.5$) around $L/D = 4$--4.5 is worth \TRADELEN{} of
apogee, about \SI{0.4}{kg} of mass margin at \SI{-143}{m/kg}. It also costs about \SI{60}{mm} of
nose-bay length (the extra length is partly recovered because more of the 3U fits inside the
curved part) and about \SI{30}{g} of shell. With the LV capped at \SI{5.4}{m}
(2027\_C\_SE\_REQ\_04) and dry-mass growth expected (the LVS flags its own mass as
underestimated; 2027\_C\_SE\_REQ\_06 caps it at \SI{50}{kg}), the design point is
\textbf{$L/D = 4.5$, $L = \SI{810}{mm}$}, which gives \TRADELEN{} of margin for \SI{62}{mm} of
bay length. If systems engineering cannot release those \SI{62}{mm}, the same profile at
$L/D = 4$ (\SI{720}{mm}, the LVS length) is the fallback; nothing else in the design changes.

\subsection{Tip}

A spherical tip of radius $r_n = \SI{4.5}{mm}$ ($r_n/R = 0.05$), machined in Al 6082-T6 as a
replaceable insert:
\begin{itemize}[leftmargin=*,itemsep=2pt]
\item it removes the sharp apex that cannot be laid up in prepreg and chips on landing;
\item it costs \BLUNTCOST{} in wave drag (\BLUNTAPOA{} of apogee). A \SI{9}{mm} tip would lower the
stagnation heating by $\sqrt2$ but cost \BLUNTAPOB{}, and the aluminium tip does not need it;
\item it stays below \TIPAL{} in the envelope case, whereas a printed PET-CF tip reaches
\PETENV{};
\item it adds about \SI{60}{g} at the very front, which helps the static margin.
\end{itemize}
